\section{Introduction}
\label{sec:introduction}

Large language models (LLMs) are being incorporated into machine learning
systems at a rapid pace: as embedding providers for retrieval and ranking,
as feature-enrichment components that turn raw text into structured
signals, and increasingly as reasoning components sitting directly in a
production decision pipeline \cite{llm4rec-survey}. Recommender-systems
research has a documented history of gains that shrink or vanish under
controlled comparison: reproducibility studies found that many reported
neural improvements did not survive against properly tuned simple
baselines \cite{dacrema-progress, rendle-ncf-mf}, that conclusions shift
with splitting, sampling, and tuning choices \cite{sun-rigorous}, and that
sampled-negative metrics can mis-rank methods \cite{krichene-sampled}.
Cost is rarely part of the comparison at all. LLM components inherit
all of these risks and add a new one: a per-item inference cost tens of
times that of the component they replace.

In practice, the behavior of LLM-augmented pipelines is far less
consistent than the enthusiasm around them would suggest. The same
architecture that provides a clear win on one dataset can be a wash, or
even a regression, on another. Practitioners deciding whether to add an
LLM component to an existing ML pipeline need more than a single
accuracy number: they need to know whether the gain is stable across
random seeds and data splits, whether it survives distribution shift, and
whether it is worth the additional latency and cost relative to a much
cheaper classical or small-model baseline.

This paper studies LLMs as \emph{components inside} ML pipelines rather
than as end-to-end systems, and asks a narrower, more actionable question
than ``do LLMs help'': \emph{when} do they help, by how much, how
reliably, and at what cost, relative to the alternatives a practitioner
would otherwise reach for.

\subsection{Contributions}
\begin{itemize}
  \item A controlled comparison of six frozen item encoders spanning
    classical TF-IDF, a 22M-parameter sentence encoder, a 335M-parameter
    embedding encoder, a 3.8B chat LLM used through its raw hidden
    states, and a 4B LLM trained for embedding. All of them feed the
    same personalized learning-to-rank pipeline, under two scorers, two
    negative-sampling regimes, a random and a temporal split, and
    popularity and matrix-factorization reference baselines
    (\S\ref{sec:pipeline-design}--\S\ref{sec:results}).
  \item Paired per-user significance and equivalence tests with a
    pre-stated margin, so that ``no gain'' claims are distinguished from
    ``no evidence'' (\S\ref{sec:evaluation-metrics}).
  \item Three findings, replicated across three datasets. (i) A scorer
    that sees only item features reduces to a popularity model, under
    which every encoder looks alike. (ii) Once ranking is personalized,
    the step from TF-IDF to a small dense encoder is the most
    consistent gain. Larger encoders add a small gain only on long
    review text with $\sim$10K users, at tens of times the cost
    (\S\ref{sec:results}). (iii) A chat LLM's raw hidden states can be
    significantly \emph{worse} than TF-IDF, depending on text length
    and scorer, while an LLM of the same size trained for embedding is
    reliable (\S\ref{sec:failure-analysis}). Within one dataset, the
    advantage of dense over lexical encoders also tracks text redundancy
    (\S\ref{sec:redundancy-analysis}).
  \item A like-for-like latency benchmark, with the same batching, inputs,
    and hardware for every encoder, and a cost-aware checklist for
    practitioners (\S\ref{sec:quality-cost},
    \S\ref{sec:practical-guidelines}).
\end{itemize}

\textbf{A cautionary note on evaluation design.} Our first finding is
methodological: feeding item embeddings to a ranker that sees nothing
about the user makes every encoder look alike, and sampling easy
negatives hides the gap between lexical and dense encoders. Either one
alone would have led us to a wrong conclusion
(\S\ref{sec:movielens-results}).

\textbf{Reproducibility.} Code, experiment scripts, per-user results,
and the scripts that generate every table are available at
\ifanon
\url{https://anonymous.4open.science/r/llm-augmented-ml-pipelines-FC52/}.
\else
\url{CAMERA-READY-CODE-LINK}.
\fi
